\section*{Capítulo \ref{ch:g9:climate-past-and-present} --- O clima da Terra, passado e presente}

\begin{solution}{exo:g9:climate-past-and-present:1}
(a) Cerca de \qty{-19}{\degreeCelsius}. (b) Cerca de
\qty{15}{\degreeCelsius}. (c) $15 - (-19) = 34$: cerca de
\qty{34}{\degreeCelsius}. (d) Dois entre: gás carbônico, vapor de água,
metano.
\end{solution}

\begin{solution}{exo:g9:climate-past-and-present:2}
(a) \omterm{def:g1:watching-the-weather:weather}{Tempo}. (b) \omterm{def:g5:weather-and-climate:climate}{Clima}. (c) \omterm{def:g5:weather-and-climate:climate}{Clima} (e \omterm{def:g9:climate-past-and-present:climate-change}{mudança climática}: uma mudança ao
longo de décadas). (d) \omterm{def:g1:watching-the-weather:weather}{Tempo}.
\end{solution}

\begin{solution}{exo:g9:climate-past-and-present:3}
(a) \qty{19}{cm}, onze anos. (b) $19 \div 11 \approx \qty{1.7}{cm}$.
(c) A \omterm{def:g4:clouds-rain-snow:precipitation}{neve} de cada ano cai por cima da \omterm{def:g4:clouds-rain-snow:precipitation}{neve} dos anos anteriores, então,
quanto mais funda uma camada, mais cedo ela caiu.
\end{solution}

\begin{solution}{exo:g9:climate-past-and-present:4}
(a) $22.4 - 21.0 = +\qty{1.4}{\degreeCelsius}$. (b) $20.3 - 21.0 =
\qty{-0.7}{\degreeCelsius}$. (c) \qty{0}{\degreeCelsius}. (d) Que esse mês
foi mais frio do que a sua média de 1991--2020.
\end{solution}

\begin{solution}{exo:g9:climate-past-and-present:5}
(a) Dos anos 1880 aos anos 1930, os anos 1950 e os anos 1960. (b) Cerca
de \qty{0.25}{\degreeCelsius} e \qty{0.80}{\degreeCelsius}. (c) $0.80 -
0.25 = \qty{0.55}{\degreeCelsius}$. (d) Não: as décadas dos anos 1880 aos
anos 1910 ficaram um pouco mais frias, e os anos 1950 foram um pouco mais
frios do que os anos 1940; mas, desde os anos 1970, cada década foi mais
quente do que a anterior.
\end{solution}

\begin{solution}{exo:g9:climate-past-and-present:6}
(a) $(324.62 - 315.98) \div 10 \approx \qty{0.86}{ppm}$ por ano. (b)
$(427.35 - 401.01) \div 10 \approx \qty{2.63}{ppm}$ por ano. (c) $2.63
\div 0.86 \approx 3.1$: cerca de três vezes mais rápida. (d) $(427.35 - 315.98)
\div 315.98 \times 100 \approx 35~\%$.
\end{solution}

\begin{solution}{exo:g9:climate-past-and-present:7}
(a) $8 \times 10^5$ anos. (b) Há cerca de \num{120000} a \num{130000}
anos: a curva da temperatura sobe acima de 0 ali. (c) Cerca de
\qty{275}{ppm}. (d) Sim: os picos e os \omterm{def:g1:landscapes-around-us:landform}{vales} das duas curvas aparecem nas
mesmas épocas.
\end{solution}

\begin{solution}{exo:g9:climate-past-and-present:8}
(a) $3.2 \times 10 = \qty{32}{mm} = \qty{3.2}{cm}$. (b) $2100 - 2025 =
75$ anos; $75 \times 3.2 = \qty{240}{mm}$, cerca de \qty{24}{cm} (e mais,
se a subida se acelerar). (c) A água do \omterm{def:g2:water-on-earth:ocean}{mar} se expande quando se aquece;
o gelo em terra derrete e a sua água escorre para o \omterm{def:g2:water-on-earth:ocean}{mar}. (d) Não: o \omterm{def:g8:ice-on-earth:sea-ice}{gelo
marinho} já flutua no \omterm{def:g2:water-on-earth:ocean}{mar}, então o seu derretimento não o faz subir.
\end{solution}

\begin{solution}{exo:g9:climate-past-and-present:9}
(a) 90. (b) 66. (c) 99. (d) \emph{Provável} quer dizer pelo menos 66
chances em 100, não certeza: o jornal faz a afirmação parecer mais segura
do que os cientistas disseram.
\end{solution}

\begin{solution}{exo:g9:climate-past-and-present:10}
Um inverno frio num lugar é \emph{\omterm{def:g1:watching-the-weather:weather}{tempo}}: diz o que aconteceu durante
alguns meses num ponto. O \omterm{def:g9:climate-past-and-present:global-warming}{aquecimento global} é uma mudança do
\emph{\omterm{def:g5:weather-and-climate:climate}{clima}}, medida ao longo de décadas, e \emph{global}: é a média
sobre a Terra inteira. Um inverno frio aqui pode acontecer enquanto a
Terra como um todo está mais quente do que antes, e as médias das
décadas continuam subindo.
\end{solution}

\begin{solution}{exo:g9:climate-past-and-present:11}
(a) $280 - 190 = \qty{90}{ppm}$, em cerca de $2 \times 10^4$ anos. (b) $427
- 280 = \qty{147}{ppm}$, em $2025 - 1750 = 275$ anos. (c) A segunda
subida é maior (147 contra 90~ppm). Ritmos: $90 \div (2 \times 10^4) =
4.5 \times 10^{-3}$~ppm por ano, contra $147 \div 275 \approx 5.3 \times
10^{-1}$~ppm por ano: cerca de 120 vezes mais rápido. (d) Que, em 2019,
ele era mais alto do que em qualquer momento de pelo menos dois milhões
de anos (alta confiança).
\end{solution}

\begin{solution}{exo:g9:climate-past-and-present:12}
(a) São quase iguais: 1.07 contra \qty{1.06}{\degreeCelsius}. A melhor
estimativa diz que a influência humana explica praticamente todo o
aquecimento observado. (b) Não: as causas naturais, juntas, mudaram a
temperatura em \qty{0.1}{\degreeCelsius} no máximo, dez vezes menos do
que o aquecimento observado. (c) Pelo menos 66 chances em 100 de que o
aquecimento causado pelo ser humano esteja entre 0.8 e
\qty{1.3}{\degreeCelsius}.
\end{solution}

\begin{solution}{pb:g9:climate-past-and-present:1}
\textbf{1.} Porque é o período de referência: as anomalias são medidas a
partir da sua própria média.

\textbf{2.} Mais frio, em \qty{0.22}{\degreeCelsius}.

\textbf{3.} Mais quente, em \qty{0.98}{\degreeCelsius}.

\textbf{4.} $0.98 - (-0.22) = \qty{1.20}{\degreeCelsius}$.

\textbf{5.} $427 - 316 = \qty{111}{ppm}$.

\textbf{6.} $111 \div 316 \times 100 \approx 35~\%$.

\textbf{7.} $427 - 300 = \qty{127}{ppm}$.

\textbf{8.} O gás carbônico absorve parte da luz infravermelha emitida
pela superfície aquecida e manda parte dela de volta para baixo. Com mais
gás carbônico no ar, mais energia volta para a superfície, que se aquece.

\textbf{9.} $2025 - 1993 = 32$ anos; $32 \times 3.2 \approx
\qty{102}{mm}$, cerca de \qty{10}{cm}.

\textbf{10.} A água do \omterm{def:g2:water-on-earth:ocean}{mar} que se aquece se expande; o gelo em terra
(\omterm{def:g8:ice-on-earth:glacier}{geleiras}, \omterm{def:g8:ice-on-earth:ice-sheet}{mantos de gelo}) derrete e escorre para o \omterm{def:g2:water-on-earth:ocean}{mar}.

\textbf{11.} Sim, do mesmo tamanho: cerca de \qty{1.2}{\degreeCelsius}
contra \qty{1.09}{\degreeCelsius}. Eles diferem um pouco porque os
períodos não são os mesmos (1880--1899 em vez de 1850--1900; 2015--2024
em vez de 2011--2020: os anos mais recentes foram mais quentes), e porque
o IPCC combina vários registros, não só o da NASA.

\textbf{12.} ``Desde o fim do século~XIX, a superfície da Terra se
aqueceu cerca de \textbf{\qty{1.2}{\degreeCelsius}}.''
\end{solution}
